\chapter{The Ledgers}
% OUTLINE Ch5 "The Ledgers" -- the series' LAST chapter. Citation ledger (Vol 5,
% per chapter; Kodo's second eyes) + THE FENCE LEDGER (this volume's own, the
% series-closing instrument: every fence held across the FIVE volumes, each row
% traced to WHERE held + WHAT it holds back -- Kodo's trace-gate: every row
% traces to a fence actually held, no ledgered fence the prose didn't hold, no
% held fence the ledger omits; Beastmaster arrival-reads it). This is the table
% the external reviewers had to reconstruct -- delivered. Gate: Kodo (trace-gate)
% + Beastmaster (arrival-read) + whole-book pass.

The earlier volumes each closed with a ledger, because a series built on
the discipline of the open receipt owes its reader an audit at the end of
every book. This volume closes with two, and the second of them is the one
the whole series has been waiting to be able to write. The first is the
citation ledger, as always: every claim in these pages traced to a name and
a date. The second is new, and it is the fitting last instrument of a
series whose spine has been its boundaries --- the fence ledger, in which
every fence the five volumes hold is set down in one place, with where it
is held and what it holds back. A series that has told the reader, book by
book, exactly where it stops, finishes by tabulating the stops.

\section*{The citation ledger}

\subsection*{Preface}
\begin{itemize}
\item David Hilbert, 1900 (International Congress of Mathematicians, Paris)
--- the sixth problem, ``the mathematical treatment of the axioms of
physics,'' named as the program this series is one finite instance of.
\end{itemize}

\subsection*{Chapter 1 --- The Single Invariant}
\begin{itemize}
\item The earned chain of the earlier volumes --- the counted slip, its
floor-and-one, the quadratic surd, the walk's stopping walls, the
reciprocal-at-scale --- gathered and cited as apparatus, not re-derived.
\item The laboratory value near $137.036$ named, once, as the world's own
measured quantity --- attributed to the laboratory and fenced from the
model's.
\end{itemize}

\subsection*{Chapter 2 --- The Finite Gauge}
\begin{itemize}
\item The comparing-bridge discipline of the first volume (Christie's
diamond, perfected by Wheatstone) --- cited as apparatus.
\item The totality property proved and named in the second volume --- the
terminating, total fragment --- cited as apparatus.
\end{itemize}

\subsection*{Chapter 3 --- The Resonances at the Fence}
\begin{itemize}
\item Claude-Louis Navier and George Gabriel Stokes --- the equations of
viscous flow; the Clay Mathematics Institute, 2000 --- the Millennium
Prize problem of their global regularity.
\item Wolfgang Pauli, 1930 --- the postulation of the neutrino; Clyde Cowan
and Frederick Reines, 1956 --- its first detection.
\item Leon Cooper, 1956 --- the electron pair; Bardeen, Cooper, and
Schrieffer, 1957 --- the theory of superconductivity built on it.
\end{itemize}

\subsection*{Chapter 4 --- The Owed and the Read}
\begin{itemize}
\item The four descriptions and the founding vow of the series' first
volume (the electron as the name of a model) --- cited as apparatus.
\end{itemize}

\section*{The fence ledger}

Every fence the five volumes hold, with the volume that holds it and the
thing it holds back. A fence is a boundary that does not shift; where a
boundary is one a later volume may cross, it is a debt, listed after.

\begin{itemize}
\item \emph{The lab-gap.} Held: Volume 1's reading, and hardest at this
volume's first chapter and again its fourth. Holds back: the model's own
inverse coupling is never the laboratory's measured value near $137.036$;
no argument bridges, approaches, or fits the one to the other.
\item \emph{The electron-model scope.} Held: the first volume's ground, and
this volume's third chapter. Holds back: the series speaks only of one
model of the electron --- not of any other particle (the neutrino among
them), force, nucleus, or cosmos.
\item \emph{The temperature.} Held: the first volume's reading and residue
chapters. Holds back: the instrument contains no temperature, and the full
physical identity of its minus-twenty-one, which that temperature would
complete, is left open.
\item \emph{Illustration, not identity.} Held: the middle three volumes'
device exhibits, and this volume's fourth chapter. Holds back: the
instrument's residue \emph{resembles} a curvature, a topological
obstruction, a caught limit, and a logical fixed point, and \emph{is}
none of them in the cited sense; the earning is owed to the device.
\item \emph{The non-funge.} Held: the fourth volume. Holds back: each
argument's residue is bound to its own representation; there is no single
residue standing behind all of them.
\item \emph{No transfer across the fluid pun.} Held: this volume's third
chapter. Holds back: the instrument's finite, decidable boundedness shares
only an English word with the open, continuous boundedness of the
Navier--Stokes problem; no theorem crosses between them.
\item \emph{Hilbert's sixth as frame, not solution.} Held: this volume's
preface. Holds back: the series is one bounded instance of the program to
axiomatize physics, and settles nothing about whether physics as a whole
can be so treated.
\item \emph{Gauge, the metrologist's word.} Held: this volume's second
chapter. Holds back: the instrument is a comparing gauge, not the
physicist's gauge --- it makes no claim in the register of field theory's
symmetries.
\item \emph{The measurement side-effect and the unresolved difference, shape
only.} Held: this volume's preface. Holds back: the instrument's
self-measurement perturbs its own count (Heisenberg's 1927 measurement
analysis) and, more rigorously, \emph{defines} its own difference by a
finite boxing under which the first and second variation are one --- a
machine-checked, identity-decided fact that resonates with Heisenberg's
non-resolvability and instantiates no quantum uncertainty. The kinship to
the fourth volume's self-fences is of family, not identity: a
finite-resolution result, not an undefinability.
\item \emph{No road to quantum gravity.} Held: this volume's third chapter.
Holds back: the series has no theory of quantum gravity and no anti-Cooper
interaction; its only artifact adjacent to the pairing idea is the one
non-terminating specimen it quarantines.
\item \emph{No teleology.} Held: the third volume's first chapter. Holds
back: an extremal principle proves a configuration extremizes a chosen
functional, never that nature prefers, intends, or deliberates.
\item \emph{No message from self-reference about minds or worlds.} Held:
the fourth volume. Holds back: the theorems of self-reference are about
formal systems, and say nothing this series will repeat about
consciousness, the limits of physics, or the soul.
\item \emph{No unexhibited choice.} Held: the first volume's preface. Holds
back: the argument admits no chooser it cannot construct; whether to accept
one is left to the reader.
\item \emph{The continuum unclaimed.} Held: the third volume's preface.
Holds back: the redundancy of ``many'' is a fact about the representation's
completeness, not a claim about the dimensionality of the world.
\end{itemize}

And the debts, the boundaries a later volume may cross, named so the fences
are not mistaken for them: the fluid description, owed to the volume that
will build the fourth face; the earning of the residue's resemblances, owed
to the device's self-founding; and the two configuring numbers and the
counted eighteen's structure, owed to the device's own deepening. A fence
is held forever; a debt is held until it is paid; the series has been
scrupulous about which is which, and this ledger keeps them apart.

\section*{The closing line}

Five volumes set out to treat one physical quantity by axioms a machine
could check, and to do it in the open --- every number earned at a named
site, every boundary drawn where honesty required, nothing claimed that was
not carried. What they built is four descriptions of one model of the
electron and a single number those descriptions converge on, read to the
model's own resolution and fenced, at the hardest point, from the world's.
What they leave is a short list of honest debts and a longer list of fences,
now tabulated, so that the reader holds the whole map of the boundaries and
does not have to draw it. The number was read. The line was not crossed.
The receipts are here, and the series is closed.
